(a) A homogeneous element $m\in M_d$ defines a global section of $\widetilde M(d)$: on $D_+(f)$, its image is $m/1\in(M(d)_f)_0$. These sections commute with multiplication by homogeneous elements of $S$, giving $\alpha:M\to\Gamma_*(\widetilde M)$.

(b) First take $M=S$ with $S$ a domain and $S_+\ne0$. Choose generators $x_1,\ldots,x_r\in S_1$. As in the proof of Theorem 5.19, the ring $S'=\bigoplus_{n\ge0}\Gamma(X,\mathcal O_X(n))$ lies in $\bigcap_iS_{x_i}$ and is a finite $S$-module. Choose homogeneous module generators of $S'$. Clearing their denominators in every $S_{x_i}$ gives an integer $N$ with $S_{\ge N}S'\subseteq S$. Since these generators have bounded degrees, $S'_d=S_d$ for all sufficiently large $d$. If $S_+=0$, both sides vanish in every positive degree.

For general $M$, take a finite graded prime filtration, with factors $(S/\mathfrak p)(a)$ for homogeneous primes $\mathfrak p$. The preceding case applies to each factor, since its associated sheaf is the direct image of the corresponding sheaf on $\operatorname{Proj}(S/\mathfrak p)$. For $0\to M'\to M\to M''\to0$, compare its degree-$d$ sequence with the left exact sequence of global sections after sheafification and twisting. If $\alpha_d$ is an isomorphism for $M'$ and $M''$, a diagram chase makes it an isomorphism for $M$. Induction proves the claim.

(c) Interpret morphisms modulo equivalence as maps of sufficiently high truncations, identified after further truncation. Sheafification is unchanged by truncation. Part (b) identifies every quasi-finitely generated $M$ with $\Gamma_*(\widetilde M)$ in this category, while Proposition 5.15 gives $\widetilde{\Gamma_*(\mathcal F)}\cong\mathcal F$.

To check that $\Gamma_*(\mathcal F)$ is quasi-finitely generated for coherent $\mathcal F$, choose a surjection from a finite sum of twists of $\mathcal O_X$ onto $\mathcal F$. Its kernel is coherent. By Theorem 5.17, it is a quotient of $\mathcal O_X(-m)^r$ for $m$ sufficiently large. By (b) for $S$, increasing $m$ ensures that every entry of the resulting map into the first sum of twists comes from $S$. Its graded cokernel is a finite graded module $M$ with $\widetilde M\cong\mathcal F$. Now apply (b). Naturality of the two identifications gives inverse equivalences.
